\documentclass[12pt,a4paper,oneside]{article}
\usepackage[brazil]{babel}
\usepackage[utf8]{inputenc}
\usepackage{olymp}
\usepackage{color}
\usepackage[dvipsnames]{xcolor}
\usepackage{listings}
\usepackage[T1]{fontenc}

\setcounter{problem}{5}

\begin{document}

\begin{problem}{Interseção}{intersecao}{2}

A interseção entre dois conjuntos $A$ e $B$, representada por $A \cap B$, é um conjunto contendo todos os elementos que pertencem, simultaneamente, aos dois conjuntos. Ou seja, um elemento $x \in A \cap B$ se e somente se $x \in A$ e $x \in B$.

Escreva um programa para calcular a interseção de dois conjuntos.

%\Notes
%
%Observações, se tiver.

\InputFile

A entrada contém três linhas: a primeira contém dois números inteiros, $N$ e $M$, representando a quantidade de elementos dos conjuntos $A$ e $B$, respectivamente; a segunda linha contém $N$ números inteiros, os elementos do conjunto $A$; e a terceira linha contém $M$ números inteiros, os elementos do conjunto $B$. Restrições: $1 \leq N, M \leq 20$.

\OutputFile

Seu programa deve gerar apenas uma linha de saída, contendo o resultado de $A \cap B$, em notação de conjuntos, isto é, envolvidos por \texttt{\{ \}}. Os elementos da interseção devem aparecer na ordem em que são apresentados no conjunto $A$, e deve haver um espaço em branco entre eles e as chaves (veja exemplos de saída). 


% \Notes
% Escreva espaço em branco para separar os números, mas não escreva espaço em branco antes do primeiro nem depois do último número da lista. Após o último escreva apenas um caractere de quebra de linha.


\Example

\begin{example}
\exmp{
5 5
4 5 3 2 1
9 4 7 1 8
}{
\{ 4 1 \}
}%
\end{example}

\begin{example}
\exmp{
5 5
89 26 37 42 11
26 11 89 42 37
}{
\{ 89 26 37 42 11 \}
}%
\end{example}

\begin{example}
\exmp{
3 5
44 22 11
20 60 85 10 2
}{
\{ \}
}%
\end{example}



%Comentários sobre o exemplo, se necessário.

\end{problem}

\end{document}
